\section{Related work}
\label{sec:related}

\subsection{gRNA efficacy models}
On-target activity prediction began with position-specific sequence models.
Rule Set 1 (Doench et al.\ 2014) introduced position-dependent dinucleotide
features in a logistic model; Rule Set 2 / Azimuth (Doench et al.\ 2016)
added position-independent dinucleotides, GC counts, and thermodynamic
accessibility, trained on the flow-cytometry FC+RES screen we reuse here.
Chari et al.\ (2015) built a support-vector model on a smaller
lentiviral screen. Deep learning entered with DeepHF (Wang et al.\ 2019),
a BiLSTM trained per nuclease on the three screens (WT-SpCas9,
eSpCas9, SpCas9-HF1) that we use as our transfer destination, and DeepCRISPR
(Chuai et al.\ 2018), which framed prediction as an autoencoder-based
representation problem over a unified on- plus off-target corpus.
CRISPRon (Kim et al.\ 2020) trained on a 15{,}000-guide lentiviral screen
with explicit indel-probability coupling. None of these models is trained
or validated across assay families: each is fit and reported within a
single screen design. Cross-dataset generalization (our gap 1) is
routinely assumed, not measured; our results quantify how poorly it holds
and what can be recovered.

\subsection{Off-target scoring}
The MIT/Hsu score (Hsu et al.\ 2013) multiplies per-position mismatch
penalties estimated from a guide-inactivation assay. The CFD score
(Doench et al.\ 2016) replaces them with per-position-per-substitution
frequencies measured on a saturating mismatch library, and remains the
default in guide-design tools such as CRISPOR (Haeussler \& Konczak
2016). Elevation (Listgarten et al.\ 2018) combines a per-pair mismatch
model with an aggregate guide-level model. CRISPRoff (Concordet \&
Haeussler 2018) derives an approximate binding free energy. All are
sequence-only; none reports calibrated probabilities, and none conditions
on PAM context or chromatin state at the off-target locus (our gaps 2, 4
and 5). Deep-learning off-target models (DeepCRISPR; CRISPR-Net, Lin et
al.\ 2020) report high accuracy on single-study benchmarks with
near-saturated negatives; our guide-grouped evaluation on crisprSQL
(Section~\ref{sec:results-offtarget}) shows how these numbers change
under a split that prevents guide leakage.

\subsection{PAM-variant and engineered-nuclease models}
Engineered Cas9 variants relax or shift the PAM requirement: SpG accepts
NGN and SpRY is near-PAMless (Walton et al.\ 2020). Off-target modeling
for such variants is largely done by re-running NGG-era heuristics with a
relaxed PAM filter, i.e.\ by ignoring the PAM channel entirely. Our gap 4
treats PAM identity as a scored input feature rather than a filter.

\subsection{pegRNA efficacy}
Prime editing (Anzalone et al.\ 2019) introduced the pegRNA design
problem. DeepPE (Kim et al.\ 2021) and PRIDICT (Mathis et al.\ 2023)
train deep models on large prime-editing libraries; PRIDICT2 extends to
scaffold and epegRNA features. Whether Cas9-efficacy knowledge transfers
to pegRNA design has been conjectured but, to our knowledge, never
benchmarked. Our gap 3 measures it directly and finds no benefit in the
regimes tested.

\subsection{Calibration of biological predictors}
Calibrated probabilities matter whenever a score triggers a downstream
cost (off-target validation assays, clinical trial enrollment).
Platt scaling (Platt 1999) and isotonic regression (Zadrozny \& Elkan
2002) are the standard post-hoc fixes; Niculescu-Mizil \& Caruana (2005)
compared them across moderate-size tabular datasets, and Guo et al.\ (2017)
documented systematic miscalibration of modern deep networks. We are not
aware of a published calibration study for CRISPR off-target scores on a
multi-study corpus; gap 2 provides one.

\subsection{Domain shift and transfer theory}
Our discovery connects to the theory of domain adaptation (Ben-David et
al.\ 2010; Pan \& Yang 2010). Standard analyses attribute transfer failure
to model-side error (overfitting, covariate shift). The target-inherent
transfer residual (Section~\ref{sec:discovery}) is a stable,
cross-replicating component of the transfer error that lives in the
\emph{assay gap itself}: it is a property of the guide, invariant to the
nuclease screen in which it is measured, and predictable from sequence
and chromatin features (the basis of the PORTA correction).
